% !TEX root = ../main.tex
% Appendix.
% Included from main.tex via % !TEX root = ../main.tex
% Appendix.
% Included from main.tex via % !TEX root = ../main.tex
% Appendix.
% Included from main.tex via % !TEX root = ../main.tex
% Appendix.
% Included from main.tex via \input{chapters/appendix}, after \appendix.

\chapter{\texorpdfstring{$L^p$ Duality}{L^p Duality}}
Let $(X,\mu)$ be a measure space.
Let $p,q,r\in\mathbb{R}$ be \textbf{H\"{o}lder triple}, i.e. they satisfy:
\begin{equation*}
    \frac{1}{p}+\frac{1}{q}=\frac{1}{r}
\end{equation*}
We have the following \textbf{H\"{o}lder's inequality}:
\begin{theorem}\label{Holder inequality}
If $f,g$ are measurable functions on $X$, then 
\begin{equation*}
    \|fg\|_{L^r}\leq \|f\|_{L^p}\|g\|_{L^q}
\end{equation*}
for $p,q,r$: H\"{o}lder triple. In particular, if $f\in L^p$, $g\in L^q$, 
then $fg\in L^r$.
\end{theorem}

In Mathlib, we often consider $E$-valued $L^p$ functions, where $E$ is a normed space (or 
real Banach space if considering their Bochner integral), for the sake of generality. So in this appendix,
we will discuss the topic up to this level of generality to be aligned with Mathlib.

For $E,F,G$: normed spaces, consider a continuous bilinear map:
\begin{equation*}
    B: E\times F\rightarrow G
\end{equation*}
it induces an associated map:
\begin{equation*}
    H_B : L^p(X;E)\times L^q(X;F)\rightarrow L^r(X;G),\quad (f,g)\mapsto (x\mapsto B(f(x),g(x)))
\end{equation*}
for $p,q,r$: H\"{o}lder triple. One can check that it is well-defined using H\"{o}lder's inequality.

For special case when $r=1$, i.e. $\frac{1}{p}+\frac{1}{q}=1$, we call $p,q$ to be \textbf{H\"{o}lder conjugate}. 
Consider the associated map $H_B: L^p(X;E)\times L^q(X;F)\rightarrow L^1(X;G)$ of the pairing $B: E\times F\rightarrow G$,
and take integral over $L^1(X;G)$ for $G$: real Banach space, we get the following \textbf{$L^p$-pairing}:
\begin{equation*}
    \Phi_B^p : L^p(X;E)\times L^q(X;F)\rightarrow G,\quad (f,g)\mapsto \int_X B(f(x),g(x)) d\mu (x).
\end{equation*}

Now let's consider the specific case for $E,F$: $\Bbbk$-normed space and $G=\Bbbk$. Then one 
can think of $B$ as a continuous linear map $E\rightarrow F^*$ and its associated $L^p$-pairing
$\Phi_B^p$ as a continuous linear map $L^p(X;E)\rightarrow (L^q(X;F))^*$. Then we have the 
following theorem:
\begin{theorem}
If $B: E\rightarrow F^*$ is an isometry, then its corresponding $L^p$-pairing
 $\Phi_B^p:L^p(X;E)\rightarrow (L^q(X;F))^*$ is also an isometry for $1\leq p<\infty$. 
If $\mu$ is semifinite, this result holds also for $p=\infty$.
\end{theorem}
\begin{proof}
We can first decompose $\Phi_B^p$ as follow:
\begin{equation*}
\begin{tikzcd}
L^p(X;E) \arrow[rd, "L^p(B)"'] \arrow[rr, "\Phi_B^p"] &                               & (L^q(X;F))^* \\
                                                    & L^p(X;F^*) \arrow[ru, "\Phi_J^p"'] &             
\end{tikzcd}
\end{equation*}
where 
\begin{equation*}
    L^p(B) : L^p(X;E)\rightarrow L^p(X;F^*),\quad f\mapsto B\circ f
\end{equation*}
and $J:F^*\times F\rightarrow \Bbbk$ is the canonical pairing and $\Phi_J^p$ is the corresponding $L^p$-pairing. One 
can check this diagram commutes. Now it suffices to show that $\Phi_J^p$ is isometry, as if so, we will have:
\begin{equation*}
    \|\Phi_B^p (f)\|=\|\Phi_J^p\circ L^p(B)(f)\|=\|B\circ f\|=\|f\|
\end{equation*}
as both $\Phi_J^p$ and $B$ are isometries.

To show $\Phi_J^p$ is isometry, i.e. for all $h\in L^p(X;F^*)$, we have $\|\Phi_J^p(h)\|=\|h\|$. First 
consider the following lemma:
\begin{lemma}
For a continuous bilinear pairing $B:E\times F\rightarrow \Bbbk$ and its
corresponding $L^p$-pairing $\Phi_B^p$, we have $\|\Phi_B^p\|\leq \|B\|$.
\end{lemma}
\begin{proof}[Proof of Lemma]
Let $p,q$ be conjugates, we have
\begin{equation*}
    |\Phi_B^p(g)(f)|\leq \int_X |B(g(x),f(x))|d\mu(x)\leq \|B\|\int_X\|g(x)\|\|f(x)\| d\mu(x)\leq \|B\|\|g\|_p\|f\|_q
\end{equation*}
where the last step uses H\"{o}lder's inequality. Thus
\begin{equation*}
    \|\Phi_B^p\|=\sup_{\|f\|_q\leq 1, \|g\|_p\leq 1}\|\Phi_B^p(g)(f)\|\leq \|B\|.
\end{equation*}
\end{proof}
Therefore, given $h\in L^p(X;F^*)$, we have 
\begin{equation*}
    \|\Phi_J^p(h)\|\leq \|\Phi_J^p\|\|h\|\leq \|J\|\|h\|\leq \|h\|
\end{equation*}
as one can show that $\|J\|\leq 1$. So now it suffices to show that for given $h\in L^p(X;F^*)$, we 
have $\|\Phi_J^p(h)\|\geq \|h\|$. 

For $1\leq p <\infty$, we first consider the following lemma:
\begin{lemma}
Given $\varepsilon>0$, for all $\varphi\in F^*$, there exists $v\in F$ s.t. $\|v\|\leq 1$ and 
$\text{Re}\;\varphi(v)\geq (1-\varepsilon)\|\varphi\|$.
\end{lemma}
\begin{proof}[Proof of Lemma]
Note $\|\varphi\|=\sup_{\|v\|\leq 1}\|\varphi(v)\|$, so for all $\varepsilon>0$, we can always find $\|v\|\leq 1$ s.t.
\begin{equation*}
    |\varphi(v)|\geq (1-\varepsilon)\|\varphi\|
\end{equation*}
Now let's consider $\tilde{v}:=\frac{\overline{\varphi(v)}}{|\varphi(v)|}v$, then $\varphi(\tilde{v})=|\varphi(v)|$, so
\begin{equation*}
    \text{Re}\;\varphi(\tilde{v})=\text{Re}\;|\varphi(v)|=|\varphi(v)|>(1-\varepsilon)\|\varphi\|.
\end{equation*}
\end{proof}
Given $\varepsilon>0$, let's consider $h$ to be a simple function first and write $h=\sum_{i=1}^n \varphi_i \1_{A_i}$ with $A_i$ pairwise disjoint and $\varphi_i\neq 0$, so that $\mu(A_i)<\infty$.
Then for each $\varphi_i\in F^*$, pick $v_i\in F$ s.t. $\|v_i\|\leq 1$ and 
$\text{Re}\;\varphi_i(v_i)\geq (1-\varepsilon)\|\varphi_i\|$. Now consider 
\begin{equation*}
    f=\sum_{i=1}^n \|\varphi_i\|^{p-1}v_i \1_{A_i}
\end{equation*}
First, we can see that $f\in L^q$, as it is a simple function supported on a set of finite measure. Moreover,
\begin{equation*}
    \|f\|_q^q=\sum_{i=1}^n\|\varphi_i\|^{q(p-1)}\|v_i\|^q \mu(A_i)\leq \sum_{i=1}^n\|\varphi_i\|^p \mu(A_i)=\|h\|_p^p
\end{equation*}
thus $\|f\|_q \leq \|h\|_p^{p/q}=\|h\|_p^{p-1}$ (Here we skip the details for specific case of 
$p=1$, $q=\infty$, but it is not hard to check). Note that 
\begin{align*}
    |\Phi_J^p(h)(f)|&\geq \text{Re}\; \Phi_J^p(h)(f)\\
      &= \sum_i \|\varphi_i\|^{p-1}\text{Re}\;\varphi_i(v_i)\mu(A_i)\\
      &\geq (1-\varepsilon)\sum_i \|\varphi_i\|^p \mu(A_i)\\
      &= (1-\varepsilon)\|h\|_p^p.
\end{align*}
so we have
\begin{equation*}
    \|\Phi_J^p(h)\|\geq \frac{|\Phi_J^p(h)(f)|}{\|f\|_q}\geq (1-\varepsilon)\frac{\|h\|_p^p}{\|f\|_q}\geq (1-\varepsilon)\frac{\|h\|_p^p}{\|h\|_p^{p-1}}=(1-\varepsilon)\|h\|.
\end{equation*}
As it is true for all given $\varepsilon>0$, we have $\|\Phi_J^p(h)\|\geq \|h\|$ for $h$: simple function. Now let's consider
\begin{equation*}
    S:=\{g\in L^p(X;F^*):\|g\|\leq \|\Phi_J^p(g)\|\}\subset L^p(X;F^*).
\end{equation*}
It is easy to see that $S$ is closed and it contains all simple functions in $L^p$. As these simple functions are dense in 
$L^p$, $S=L^p$, which gives us desired result.

For $p=\infty$, assume $\mu$ is semifinite. We want to show that $\|\Phi_{J}^\infty(h)\|\geq \|h\|_\infty$ for 
$\Phi_{J}^\infty : L^\infty(X;F^*)\rightarrow (L^1(X;F))^*$. First we introduce the following lemma:
\begin{lemma}
For all measurable $D\subset X$ with $0<\mu(D)<\infty$, we have 
\begin{equation*}
    \|\Phi^\infty_J(h)\|\geq \frac{1}{\mu(D)}\int_D \|h(x)\| d\mu(x).
\end{equation*}
\end{lemma}
\begin{proof}[Proof of Lemma]
Given measurable $D\subset X$ with $0<\mu(D)<\infty$. Consider the operator 
\begin{equation*}
    \iota_D : L^\infty(X) \rightarrow L^1(X), \quad f\mapsto \1_D f
\end{equation*}
It is well-defined as $\mu(D)$ is finite and $f\in L^\infty$ is bounded a.e.. Then we can show that the following diagram
\begin{equation*}
\begin{tikzcd}
L^\infty (X;F^*) \arrow[rr, "\Phi_J^\infty"] \arrow[d, "\iota_D"'] &  & (L^1(X;F))^* \arrow[d, "\iota_D^*"] \\
L^1 (X;F^*) \arrow[rr, "\Phi_J^1"]                                 &  & (L^\infty(X;F))^*                  
\end{tikzcd}
\end{equation*}
commutes. This can be shown by:
\begin{align*}
  \Phi_J^1(\1_D h)(f)&=\int_X \1_D(x)h(x)(f(x))d\mu(x)\\
  &=\int_D h(x)(\1_D(x)f(x))d\mu(x)=\Phi^\infty_J(h)(\1_D f).
\end{align*}
Therefore,
\begin{equation*}
    |\Phi_J^1(\1_D h)(f)|=|\Phi^\infty_J(h)(\1_D f)|\leq \|\Phi^\infty_J(h)\|\|\1_D f\|_1\leq \mu(D)\|\Phi^\infty_J(h)\|\| f\|_\infty
\end{equation*}
This implies:
\begin{equation*}
    \|\Phi_J^1(\1_D h)\|=\sup_{\|f\|_\infty \leq 1}|\Phi_J^1(\1_D h)(f)|\leq \sup_{\|f\|_\infty \leq 1}\mu(D)\|\Phi^\infty_J(h)\|\| f\|_\infty\leq \mu(D)\|\Phi^\infty_J(h)\|.
\end{equation*}
Recall that we've already proved that 
\begin{equation*}
    \|\Phi_J^1(\1_D h)\|=\|\1_D h\|_1=\int_D \|h(x)\|d\mu(x)
\end{equation*}
Thus we have our desired result.
\end{proof}
Now let $M:=\|h\|_\infty$. For $M=0$, the inequality trivially holds. For $M>0$, given $0<\varepsilon<M$, consider
\begin{equation*}
    A:=\{x\in X: \|h(x)\|>M-\varepsilon\}
\end{equation*}
By definition of $\|\cd\|_\infty$, $\mu(A)>0$. As $\mu$ is semifinite, there exists $D\subset A$: measurable s.t. $0<\mu(D)<\infty$.
As $D\subset A$, for all $x\in D$, we have $\|h(x)\|>M-\varepsilon$. Therefore,
\begin{equation*}
    \|\Phi^\infty_J(h)\|\geq \frac{1}{\mu(D)}\int_D \|h(x)\|d\mu(x)\geq M-\varepsilon.
\end{equation*}
As it is true for all given $\varepsilon$, we have $\|\Phi_J^\infty(h)\|\geq \|h\|_\infty$.
\end{proof}


\begin{theorem}
Let $p$ and $q$ be H\"{o}lder conjugate, and let $B:\Bbbk\times\Bbbk\rightarrow\Bbbk$ be the multiplication.
If $1<p<\infty$, for each $\phi\in(L^q(X;\Bbbk))^*$ there exists $f\in L^p(X;\Bbbk)$ such that
$\phi(g)=\int_X fg\, d\mu$ for all $g\in L^q(X;\Bbbk)$, and hence
$\Phi_B^p:L^p(X;\Bbbk)\rightarrow(L^q(X;\Bbbk))^*$ is an isometric isomorphism.
The same conclusion holds for $p=\infty$ provided $\mu$ is $\sigma$-finite.
\end{theorem}















, after \appendix.

\chapter{\texorpdfstring{$L^p$ Duality}{L^p Duality}}
Let $(X,\mu)$ be a measure space.
Let $p,q,r\in\mathbb{R}$ be \textbf{H\"{o}lder triple}, i.e. they satisfy:
\begin{equation*}
    \frac{1}{p}+\frac{1}{q}=\frac{1}{r}
\end{equation*}
We have the following \textbf{H\"{o}lder's inequality}:
\begin{theorem}\label{Holder inequality}
If $f,g$ are measurable functions on $X$, then 
\begin{equation*}
    \|fg\|_{L^r}\leq \|f\|_{L^p}\|g\|_{L^q}
\end{equation*}
for $p,q,r$: H\"{o}lder triple. In particular, if $f\in L^p$, $g\in L^q$, 
then $fg\in L^r$.
\end{theorem}

In Mathlib, we often consider $E$-valued $L^p$ functions, where $E$ is a normed space (or 
real Banach space if considering their Bochner integral), for the sake of generality. So in this appendix,
we will discuss the topic up to this level of generality to be aligned with Mathlib.

For $E,F,G$: normed spaces, consider a continuous bilinear map:
\begin{equation*}
    B: E\times F\rightarrow G
\end{equation*}
it induces an associated map:
\begin{equation*}
    H_B : L^p(X;E)\times L^q(X;F)\rightarrow L^r(X;G),\quad (f,g)\mapsto (x\mapsto B(f(x),g(x)))
\end{equation*}
for $p,q,r$: H\"{o}lder triple. One can check that it is well-defined using H\"{o}lder's inequality.

For special case when $r=1$, i.e. $\frac{1}{p}+\frac{1}{q}=1$, we call $p,q$ to be \textbf{H\"{o}lder conjugate}. 
Consider the associated map $H_B: L^p(X;E)\times L^q(X;F)\rightarrow L^1(X;G)$ of the pairing $B: E\times F\rightarrow G$,
and take integral over $L^1(X;G)$ for $G$: real Banach space, we get the following \textbf{$L^p$-pairing}:
\begin{equation*}
    \Phi_B^p : L^p(X;E)\times L^q(X;F)\rightarrow G,\quad (f,g)\mapsto \int_X B(f(x),g(x)) d\mu (x).
\end{equation*}

Now let's consider the specific case for $E,F$: $\Bbbk$-normed space and $G=\Bbbk$. Then one 
can think of $B$ as a continuous linear map $E\rightarrow F^*$ and its associated $L^p$-pairing
$\Phi_B^p$ as a continuous linear map $L^p(X;E)\rightarrow (L^q(X;F))^*$. Then we have the 
following theorem:
\begin{theorem}
If $B: E\rightarrow F^*$ is an isometry, then its corresponding $L^p$-pairing
 $\Phi_B^p:L^p(X;E)\rightarrow (L^q(X;F))^*$ is also an isometry for $1\leq p<\infty$. 
If $\mu$ is semifinite, this result holds also for $p=\infty$.
\end{theorem}
\begin{proof}
We can first decompose $\Phi_B^p$ as follow:
\begin{equation*}
\begin{tikzcd}
L^p(X;E) \arrow[rd, "L^p(B)"'] \arrow[rr, "\Phi_B^p"] &                               & (L^q(X;F))^* \\
                                                    & L^p(X;F^*) \arrow[ru, "\Phi_J^p"'] &             
\end{tikzcd}
\end{equation*}
where 
\begin{equation*}
    L^p(B) : L^p(X;E)\rightarrow L^p(X;F^*),\quad f\mapsto B\circ f
\end{equation*}
and $J:F^*\times F\rightarrow \Bbbk$ is the canonical pairing and $\Phi_J^p$ is the corresponding $L^p$-pairing. One 
can check this diagram commutes. Now it suffices to show that $\Phi_J^p$ is isometry, as if so, we will have:
\begin{equation*}
    \|\Phi_B^p (f)\|=\|\Phi_J^p\circ L^p(B)(f)\|=\|B\circ f\|=\|f\|
\end{equation*}
as both $\Phi_J^p$ and $B$ are isometries.

To show $\Phi_J^p$ is isometry, i.e. for all $h\in L^p(X;F^*)$, we have $\|\Phi_J^p(h)\|=\|h\|$. First 
consider the following lemma:
\begin{lemma}
For a continuous bilinear pairing $B:E\times F\rightarrow \Bbbk$ and its
corresponding $L^p$-pairing $\Phi_B^p$, we have $\|\Phi_B^p\|\leq \|B\|$.
\end{lemma}
\begin{proof}[Proof of Lemma]
Let $p,q$ be conjugates, we have
\begin{equation*}
    |\Phi_B^p(g)(f)|\leq \int_X |B(g(x),f(x))|d\mu(x)\leq \|B\|\int_X\|g(x)\|\|f(x)\| d\mu(x)\leq \|B\|\|g\|_p\|f\|_q
\end{equation*}
where the last step uses H\"{o}lder's inequality. Thus
\begin{equation*}
    \|\Phi_B^p\|=\sup_{\|f\|_q\leq 1, \|g\|_p\leq 1}\|\Phi_B^p(g)(f)\|\leq \|B\|.
\end{equation*}
\end{proof}
Therefore, given $h\in L^p(X;F^*)$, we have 
\begin{equation*}
    \|\Phi_J^p(h)\|\leq \|\Phi_J^p\|\|h\|\leq \|J\|\|h\|\leq \|h\|
\end{equation*}
as one can show that $\|J\|\leq 1$. So now it suffices to show that for given $h\in L^p(X;F^*)$, we 
have $\|\Phi_J^p(h)\|\geq \|h\|$. 

For $1\leq p <\infty$, we first consider the following lemma:
\begin{lemma}
Given $\varepsilon>0$, for all $\varphi\in F^*$, there exists $v\in F$ s.t. $\|v\|\leq 1$ and 
$\text{Re}\;\varphi(v)\geq (1-\varepsilon)\|\varphi\|$.
\end{lemma}
\begin{proof}[Proof of Lemma]
Note $\|\varphi\|=\sup_{\|v\|\leq 1}\|\varphi(v)\|$, so for all $\varepsilon>0$, we can always find $\|v\|\leq 1$ s.t.
\begin{equation*}
    |\varphi(v)|\geq (1-\varepsilon)\|\varphi\|
\end{equation*}
Now let's consider $\tilde{v}:=\frac{\overline{\varphi(v)}}{|\varphi(v)|}v$, then $\varphi(\tilde{v})=|\varphi(v)|$, so
\begin{equation*}
    \text{Re}\;\varphi(\tilde{v})=\text{Re}\;|\varphi(v)|=|\varphi(v)|>(1-\varepsilon)\|\varphi\|.
\end{equation*}
\end{proof}
Given $\varepsilon>0$, let's consider $h$ to be a simple function first and write $h=\sum_{i=1}^n \varphi_i \1_{A_i}$ with $A_i$ pairwise disjoint and $\varphi_i\neq 0$, so that $\mu(A_i)<\infty$.
Then for each $\varphi_i\in F^*$, pick $v_i\in F$ s.t. $\|v_i\|\leq 1$ and 
$\text{Re}\;\varphi_i(v_i)\geq (1-\varepsilon)\|\varphi_i\|$. Now consider 
\begin{equation*}
    f=\sum_{i=1}^n \|\varphi_i\|^{p-1}v_i \1_{A_i}
\end{equation*}
First, we can see that $f\in L^q$, as it is a simple function supported on a set of finite measure. Moreover,
\begin{equation*}
    \|f\|_q^q=\sum_{i=1}^n\|\varphi_i\|^{q(p-1)}\|v_i\|^q \mu(A_i)\leq \sum_{i=1}^n\|\varphi_i\|^p \mu(A_i)=\|h\|_p^p
\end{equation*}
thus $\|f\|_q \leq \|h\|_p^{p/q}=\|h\|_p^{p-1}$ (Here we skip the details for specific case of 
$p=1$, $q=\infty$, but it is not hard to check). Note that 
\begin{align*}
    |\Phi_J^p(h)(f)|&\geq \text{Re}\; \Phi_J^p(h)(f)\\
      &= \sum_i \|\varphi_i\|^{p-1}\text{Re}\;\varphi_i(v_i)\mu(A_i)\\
      &\geq (1-\varepsilon)\sum_i \|\varphi_i\|^p \mu(A_i)\\
      &= (1-\varepsilon)\|h\|_p^p.
\end{align*}
so we have
\begin{equation*}
    \|\Phi_J^p(h)\|\geq \frac{|\Phi_J^p(h)(f)|}{\|f\|_q}\geq (1-\varepsilon)\frac{\|h\|_p^p}{\|f\|_q}\geq (1-\varepsilon)\frac{\|h\|_p^p}{\|h\|_p^{p-1}}=(1-\varepsilon)\|h\|.
\end{equation*}
As it is true for all given $\varepsilon>0$, we have $\|\Phi_J^p(h)\|\geq \|h\|$ for $h$: simple function. Now let's consider
\begin{equation*}
    S:=\{g\in L^p(X;F^*):\|g\|\leq \|\Phi_J^p(g)\|\}\subset L^p(X;F^*).
\end{equation*}
It is easy to see that $S$ is closed and it contains all simple functions in $L^p$. As these simple functions are dense in 
$L^p$, $S=L^p$, which gives us desired result.

For $p=\infty$, assume $\mu$ is semifinite. We want to show that $\|\Phi_{J}^\infty(h)\|\geq \|h\|_\infty$ for 
$\Phi_{J}^\infty : L^\infty(X;F^*)\rightarrow (L^1(X;F))^*$. First we introduce the following lemma:
\begin{lemma}
For all measurable $D\subset X$ with $0<\mu(D)<\infty$, we have 
\begin{equation*}
    \|\Phi^\infty_J(h)\|\geq \frac{1}{\mu(D)}\int_D \|h(x)\| d\mu(x).
\end{equation*}
\end{lemma}
\begin{proof}[Proof of Lemma]
Given measurable $D\subset X$ with $0<\mu(D)<\infty$. Consider the operator 
\begin{equation*}
    \iota_D : L^\infty(X) \rightarrow L^1(X), \quad f\mapsto \1_D f
\end{equation*}
It is well-defined as $\mu(D)$ is finite and $f\in L^\infty$ is bounded a.e.. Then we can show that the following diagram
\begin{equation*}
\begin{tikzcd}
L^\infty (X;F^*) \arrow[rr, "\Phi_J^\infty"] \arrow[d, "\iota_D"'] &  & (L^1(X;F))^* \arrow[d, "\iota_D^*"] \\
L^1 (X;F^*) \arrow[rr, "\Phi_J^1"]                                 &  & (L^\infty(X;F))^*                  
\end{tikzcd}
\end{equation*}
commutes. This can be shown by:
\begin{align*}
  \Phi_J^1(\1_D h)(f)&=\int_X \1_D(x)h(x)(f(x))d\mu(x)\\
  &=\int_D h(x)(\1_D(x)f(x))d\mu(x)=\Phi^\infty_J(h)(\1_D f).
\end{align*}
Therefore,
\begin{equation*}
    |\Phi_J^1(\1_D h)(f)|=|\Phi^\infty_J(h)(\1_D f)|\leq \|\Phi^\infty_J(h)\|\|\1_D f\|_1\leq \mu(D)\|\Phi^\infty_J(h)\|\| f\|_\infty
\end{equation*}
This implies:
\begin{equation*}
    \|\Phi_J^1(\1_D h)\|=\sup_{\|f\|_\infty \leq 1}|\Phi_J^1(\1_D h)(f)|\leq \sup_{\|f\|_\infty \leq 1}\mu(D)\|\Phi^\infty_J(h)\|\| f\|_\infty\leq \mu(D)\|\Phi^\infty_J(h)\|.
\end{equation*}
Recall that we've already proved that 
\begin{equation*}
    \|\Phi_J^1(\1_D h)\|=\|\1_D h\|_1=\int_D \|h(x)\|d\mu(x)
\end{equation*}
Thus we have our desired result.
\end{proof}
Now let $M:=\|h\|_\infty$. For $M=0$, the inequality trivially holds. For $M>0$, given $0<\varepsilon<M$, consider
\begin{equation*}
    A:=\{x\in X: \|h(x)\|>M-\varepsilon\}
\end{equation*}
By definition of $\|\cd\|_\infty$, $\mu(A)>0$. As $\mu$ is semifinite, there exists $D\subset A$: measurable s.t. $0<\mu(D)<\infty$.
As $D\subset A$, for all $x\in D$, we have $\|h(x)\|>M-\varepsilon$. Therefore,
\begin{equation*}
    \|\Phi^\infty_J(h)\|\geq \frac{1}{\mu(D)}\int_D \|h(x)\|d\mu(x)\geq M-\varepsilon.
\end{equation*}
As it is true for all given $\varepsilon$, we have $\|\Phi_J^\infty(h)\|\geq \|h\|_\infty$.
\end{proof}


\begin{theorem}
Let $p$ and $q$ be H\"{o}lder conjugate, and let $B:\Bbbk\times\Bbbk\rightarrow\Bbbk$ be the multiplication.
If $1<p<\infty$, for each $\phi\in(L^q(X;\Bbbk))^*$ there exists $f\in L^p(X;\Bbbk)$ such that
$\phi(g)=\int_X fg\, d\mu$ for all $g\in L^q(X;\Bbbk)$, and hence
$\Phi_B^p:L^p(X;\Bbbk)\rightarrow(L^q(X;\Bbbk))^*$ is an isometric isomorphism.
The same conclusion holds for $p=\infty$ provided $\mu$ is $\sigma$-finite.
\end{theorem}















, after \appendix.

\chapter{\texorpdfstring{$L^p$ Duality}{L^p Duality}}
Let $(X,\mu)$ be a measure space.
Let $p,q,r\in\mathbb{R}$ be \textbf{H\"{o}lder triple}, i.e. they satisfy:
\begin{equation*}
    \frac{1}{p}+\frac{1}{q}=\frac{1}{r}
\end{equation*}
We have the following \textbf{H\"{o}lder's inequality}:
\begin{theorem}\label{Holder inequality}
If $f,g$ are measurable functions on $X$, then 
\begin{equation*}
    \|fg\|_{L^r}\leq \|f\|_{L^p}\|g\|_{L^q}
\end{equation*}
for $p,q,r$: H\"{o}lder triple. In particular, if $f\in L^p$, $g\in L^q$, 
then $fg\in L^r$.
\end{theorem}

In Mathlib, we often consider $E$-valued $L^p$ functions, where $E$ is a normed space (or 
real Banach space if considering their Bochner integral), for the sake of generality. So in this appendix,
we will discuss the topic up to this level of generality to be aligned with Mathlib.

For $E,F,G$: normed spaces, consider a continuous bilinear map:
\begin{equation*}
    B: E\times F\rightarrow G
\end{equation*}
it induces an associated map:
\begin{equation*}
    H_B : L^p(X;E)\times L^q(X;F)\rightarrow L^r(X;G),\quad (f,g)\mapsto (x\mapsto B(f(x),g(x)))
\end{equation*}
for $p,q,r$: H\"{o}lder triple. One can check that it is well-defined using H\"{o}lder's inequality.

For special case when $r=1$, i.e. $\frac{1}{p}+\frac{1}{q}=1$, we call $p,q$ to be \textbf{H\"{o}lder conjugate}. 
Consider the associated map $H_B: L^p(X;E)\times L^q(X;F)\rightarrow L^1(X;G)$ of the pairing $B: E\times F\rightarrow G$,
and take integral over $L^1(X;G)$ for $G$: real Banach space, we get the following \textbf{$L^p$-pairing}:
\begin{equation*}
    \Phi_B^p : L^p(X;E)\times L^q(X;F)\rightarrow G,\quad (f,g)\mapsto \int_X B(f(x),g(x)) d\mu (x).
\end{equation*}

Now let's consider the specific case for $E,F$: $\Bbbk$-normed space and $G=\Bbbk$. Then one 
can think of $B$ as a continuous linear map $E\rightarrow F^*$ and its associated $L^p$-pairing
$\Phi_B^p$ as a continuous linear map $L^p(X;E)\rightarrow (L^q(X;F))^*$. Then we have the 
following theorem:
\begin{theorem}
If $B: E\rightarrow F^*$ is an isometry, then its corresponding $L^p$-pairing
 $\Phi_B^p:L^p(X;E)\rightarrow (L^q(X;F))^*$ is also an isometry for $1\leq p<\infty$. 
If $\mu$ is semifinite, this result holds also for $p=\infty$.
\end{theorem}
\begin{proof}
We can first decompose $\Phi_B^p$ as follow:
\begin{equation*}
\begin{tikzcd}
L^p(X;E) \arrow[rd, "L^p(B)"'] \arrow[rr, "\Phi_B^p"] &                               & (L^q(X;F))^* \\
                                                    & L^p(X;F^*) \arrow[ru, "\Phi_J^p"'] &             
\end{tikzcd}
\end{equation*}
where 
\begin{equation*}
    L^p(B) : L^p(X;E)\rightarrow L^p(X;F^*),\quad f\mapsto B\circ f
\end{equation*}
and $J:F^*\times F\rightarrow \Bbbk$ is the canonical pairing and $\Phi_J^p$ is the corresponding $L^p$-pairing. One 
can check this diagram commutes. Now it suffices to show that $\Phi_J^p$ is isometry, as if so, we will have:
\begin{equation*}
    \|\Phi_B^p (f)\|=\|\Phi_J^p\circ L^p(B)(f)\|=\|B\circ f\|=\|f\|
\end{equation*}
as both $\Phi_J^p$ and $B$ are isometries.

To show $\Phi_J^p$ is isometry, i.e. for all $h\in L^p(X;F^*)$, we have $\|\Phi_J^p(h)\|=\|h\|$. First 
consider the following lemma:
\begin{lemma}
For a continuous bilinear pairing $B:E\times F\rightarrow \Bbbk$ and its
corresponding $L^p$-pairing $\Phi_B^p$, we have $\|\Phi_B^p\|\leq \|B\|$.
\end{lemma}
\begin{proof}[Proof of Lemma]
Let $p,q$ be conjugates, we have
\begin{equation*}
    |\Phi_B^p(g)(f)|\leq \int_X |B(g(x),f(x))|d\mu(x)\leq \|B\|\int_X\|g(x)\|\|f(x)\| d\mu(x)\leq \|B\|\|g\|_p\|f\|_q
\end{equation*}
where the last step uses H\"{o}lder's inequality. Thus
\begin{equation*}
    \|\Phi_B^p\|=\sup_{\|f\|_q\leq 1, \|g\|_p\leq 1}\|\Phi_B^p(g)(f)\|\leq \|B\|.
\end{equation*}
\end{proof}
Therefore, given $h\in L^p(X;F^*)$, we have 
\begin{equation*}
    \|\Phi_J^p(h)\|\leq \|\Phi_J^p\|\|h\|\leq \|J\|\|h\|\leq \|h\|
\end{equation*}
as one can show that $\|J\|\leq 1$. So now it suffices to show that for given $h\in L^p(X;F^*)$, we 
have $\|\Phi_J^p(h)\|\geq \|h\|$. 

For $1\leq p <\infty$, we first consider the following lemma:
\begin{lemma}
Given $\varepsilon>0$, for all $\varphi\in F^*$, there exists $v\in F$ s.t. $\|v\|\leq 1$ and 
$\text{Re}\;\varphi(v)\geq (1-\varepsilon)\|\varphi\|$.
\end{lemma}
\begin{proof}[Proof of Lemma]
Note $\|\varphi\|=\sup_{\|v\|\leq 1}\|\varphi(v)\|$, so for all $\varepsilon>0$, we can always find $\|v\|\leq 1$ s.t.
\begin{equation*}
    |\varphi(v)|\geq (1-\varepsilon)\|\varphi\|
\end{equation*}
Now let's consider $\tilde{v}:=\frac{\overline{\varphi(v)}}{|\varphi(v)|}v$, then $\varphi(\tilde{v})=|\varphi(v)|$, so
\begin{equation*}
    \text{Re}\;\varphi(\tilde{v})=\text{Re}\;|\varphi(v)|=|\varphi(v)|>(1-\varepsilon)\|\varphi\|.
\end{equation*}
\end{proof}
Given $\varepsilon>0$, let's consider $h$ to be a simple function first and write $h=\sum_{i=1}^n \varphi_i \1_{A_i}$ with $A_i$ pairwise disjoint and $\varphi_i\neq 0$, so that $\mu(A_i)<\infty$.
Then for each $\varphi_i\in F^*$, pick $v_i\in F$ s.t. $\|v_i\|\leq 1$ and 
$\text{Re}\;\varphi_i(v_i)\geq (1-\varepsilon)\|\varphi_i\|$. Now consider 
\begin{equation*}
    f=\sum_{i=1}^n \|\varphi_i\|^{p-1}v_i \1_{A_i}
\end{equation*}
First, we can see that $f\in L^q$, as it is a simple function supported on a set of finite measure. Moreover,
\begin{equation*}
    \|f\|_q^q=\sum_{i=1}^n\|\varphi_i\|^{q(p-1)}\|v_i\|^q \mu(A_i)\leq \sum_{i=1}^n\|\varphi_i\|^p \mu(A_i)=\|h\|_p^p
\end{equation*}
thus $\|f\|_q \leq \|h\|_p^{p/q}=\|h\|_p^{p-1}$ (Here we skip the details for specific case of 
$p=1$, $q=\infty$, but it is not hard to check). Note that 
\begin{align*}
    |\Phi_J^p(h)(f)|&\geq \text{Re}\; \Phi_J^p(h)(f)\\
      &= \sum_i \|\varphi_i\|^{p-1}\text{Re}\;\varphi_i(v_i)\mu(A_i)\\
      &\geq (1-\varepsilon)\sum_i \|\varphi_i\|^p \mu(A_i)\\
      &= (1-\varepsilon)\|h\|_p^p.
\end{align*}
so we have
\begin{equation*}
    \|\Phi_J^p(h)\|\geq \frac{|\Phi_J^p(h)(f)|}{\|f\|_q}\geq (1-\varepsilon)\frac{\|h\|_p^p}{\|f\|_q}\geq (1-\varepsilon)\frac{\|h\|_p^p}{\|h\|_p^{p-1}}=(1-\varepsilon)\|h\|.
\end{equation*}
As it is true for all given $\varepsilon>0$, we have $\|\Phi_J^p(h)\|\geq \|h\|$ for $h$: simple function. Now let's consider
\begin{equation*}
    S:=\{g\in L^p(X;F^*):\|g\|\leq \|\Phi_J^p(g)\|\}\subset L^p(X;F^*).
\end{equation*}
It is easy to see that $S$ is closed and it contains all simple functions in $L^p$. As these simple functions are dense in 
$L^p$, $S=L^p$, which gives us desired result.

For $p=\infty$, assume $\mu$ is semifinite. We want to show that $\|\Phi_{J}^\infty(h)\|\geq \|h\|_\infty$ for 
$\Phi_{J}^\infty : L^\infty(X;F^*)\rightarrow (L^1(X;F))^*$. First we introduce the following lemma:
\begin{lemma}
For all measurable $D\subset X$ with $0<\mu(D)<\infty$, we have 
\begin{equation*}
    \|\Phi^\infty_J(h)\|\geq \frac{1}{\mu(D)}\int_D \|h(x)\| d\mu(x).
\end{equation*}
\end{lemma}
\begin{proof}[Proof of Lemma]
Given measurable $D\subset X$ with $0<\mu(D)<\infty$. Consider the operator 
\begin{equation*}
    \iota_D : L^\infty(X) \rightarrow L^1(X), \quad f\mapsto \1_D f
\end{equation*}
It is well-defined as $\mu(D)$ is finite and $f\in L^\infty$ is bounded a.e.. Then we can show that the following diagram
\begin{equation*}
\begin{tikzcd}
L^\infty (X;F^*) \arrow[rr, "\Phi_J^\infty"] \arrow[d, "\iota_D"'] &  & (L^1(X;F))^* \arrow[d, "\iota_D^*"] \\
L^1 (X;F^*) \arrow[rr, "\Phi_J^1"]                                 &  & (L^\infty(X;F))^*                  
\end{tikzcd}
\end{equation*}
commutes. This can be shown by:
\begin{align*}
  \Phi_J^1(\1_D h)(f)&=\int_X \1_D(x)h(x)(f(x))d\mu(x)\\
  &=\int_D h(x)(\1_D(x)f(x))d\mu(x)=\Phi^\infty_J(h)(\1_D f).
\end{align*}
Therefore,
\begin{equation*}
    |\Phi_J^1(\1_D h)(f)|=|\Phi^\infty_J(h)(\1_D f)|\leq \|\Phi^\infty_J(h)\|\|\1_D f\|_1\leq \mu(D)\|\Phi^\infty_J(h)\|\| f\|_\infty
\end{equation*}
This implies:
\begin{equation*}
    \|\Phi_J^1(\1_D h)\|=\sup_{\|f\|_\infty \leq 1}|\Phi_J^1(\1_D h)(f)|\leq \sup_{\|f\|_\infty \leq 1}\mu(D)\|\Phi^\infty_J(h)\|\| f\|_\infty\leq \mu(D)\|\Phi^\infty_J(h)\|.
\end{equation*}
Recall that we've already proved that 
\begin{equation*}
    \|\Phi_J^1(\1_D h)\|=\|\1_D h\|_1=\int_D \|h(x)\|d\mu(x)
\end{equation*}
Thus we have our desired result.
\end{proof}
Now let $M:=\|h\|_\infty$. For $M=0$, the inequality trivially holds. For $M>0$, given $0<\varepsilon<M$, consider
\begin{equation*}
    A:=\{x\in X: \|h(x)\|>M-\varepsilon\}
\end{equation*}
By definition of $\|\cd\|_\infty$, $\mu(A)>0$. As $\mu$ is semifinite, there exists $D\subset A$: measurable s.t. $0<\mu(D)<\infty$.
As $D\subset A$, for all $x\in D$, we have $\|h(x)\|>M-\varepsilon$. Therefore,
\begin{equation*}
    \|\Phi^\infty_J(h)\|\geq \frac{1}{\mu(D)}\int_D \|h(x)\|d\mu(x)\geq M-\varepsilon.
\end{equation*}
As it is true for all given $\varepsilon$, we have $\|\Phi_J^\infty(h)\|\geq \|h\|_\infty$.
\end{proof}


\begin{theorem}
Let $p$ and $q$ be H\"{o}lder conjugate, and let $B:\Bbbk\times\Bbbk\rightarrow\Bbbk$ be the multiplication.
If $1<p<\infty$, for each $\phi\in(L^q(X;\Bbbk))^*$ there exists $f\in L^p(X;\Bbbk)$ such that
$\phi(g)=\int_X fg\, d\mu$ for all $g\in L^q(X;\Bbbk)$, and hence
$\Phi_B^p:L^p(X;\Bbbk)\rightarrow(L^q(X;\Bbbk))^*$ is an isometric isomorphism.
The same conclusion holds for $p=\infty$ provided $\mu$ is $\sigma$-finite.
\end{theorem}















, after \appendix.

\chapter{\texorpdfstring{$L^p$ Duality}{L^p Duality}}
Let $(X,\mu)$ be a measure space.
Let $p,q,r\in\mathbb{R}$ be \textbf{H\"{o}lder triple}, i.e. they satisfy:
\begin{equation*}
    \frac{1}{p}+\frac{1}{q}=\frac{1}{r}
\end{equation*}
We have the following \textbf{H\"{o}lder's inequality}:
\begin{theorem}\label{Holder inequality}
If $f,g$ are measurable functions on $X$, then 
\begin{equation*}
    \|fg\|_{L^r}\leq \|f\|_{L^p}\|g\|_{L^q}
\end{equation*}
for $p,q,r$: H\"{o}lder triple. In particular, if $f\in L^p$, $g\in L^q$, 
then $fg\in L^r$.
\end{theorem}

In Mathlib, we often consider $E$-valued $L^p$ functions, where $E$ is a normed space (or 
real Banach space if considering their Bochner integral), for the sake of generality. So in this appendix,
we will discuss the topic up to this level of generality to be aligned with Mathlib.

For $E,F,G$: normed spaces, consider a continuous bilinear map:
\begin{equation*}
    B: E\times F\rightarrow G
\end{equation*}
it induces an associated map:
\begin{equation*}
    H_B : L^p(X;E)\times L^q(X;F)\rightarrow L^r(X;G),\quad (f,g)\mapsto (x\mapsto B(f(x),g(x)))
\end{equation*}
for $p,q,r$: H\"{o}lder triple. One can check that it is well-defined using H\"{o}lder's inequality.

For special case when $r=1$, i.e. $\frac{1}{p}+\frac{1}{q}=1$, we call $p,q$ to be \textbf{H\"{o}lder conjugate}. 
Consider the associated map $H_B: L^p(X;E)\times L^q(X;F)\rightarrow L^1(X;G)$ of the pairing $B: E\times F\rightarrow G$,
and take integral over $L^1(X;G)$ for $G$: real Banach space, we get the following \textbf{$L^p$-pairing}:
\begin{equation*}
    \Phi_B^p : L^p(X;E)\times L^q(X;F)\rightarrow G,\quad (f,g)\mapsto \int_X B(f(x),g(x)) d\mu (x).
\end{equation*}

Now let's consider the specific case for $E,F$: $\Bbbk$-normed space and $G=\Bbbk$. Then one 
can think of $B$ as a continuous linear map $E\rightarrow F^*$ and its associated $L^p$-pairing
$\Phi_B^p$ as a continuous linear map $L^p(X;E)\rightarrow (L^q(X;F))^*$. Then we have the 
following theorem:
\begin{theorem}
If $B: E\rightarrow F^*$ is an isometry, then its corresponding $L^p$-pairing
 $\Phi_B^p:L^p(X;E)\rightarrow (L^q(X;F))^*$ is also an isometry for $1\leq p<\infty$. 
If $\mu$ is semifinite, this result holds also for $p=\infty$.
\end{theorem}
\begin{proof}
We can first decompose $\Phi_B^p$ as follow:
\begin{equation*}
\begin{tikzcd}
L^p(X;E) \arrow[rd, "L^p(B)"'] \arrow[rr, "\Phi_B^p"] &                               & (L^q(X;F))^* \\
                                                    & L^p(X;F^*) \arrow[ru, "\Phi_J^p"'] &             
\end{tikzcd}
\end{equation*}
where 
\begin{equation*}
    L^p(B) : L^p(X;E)\rightarrow L^p(X;F^*),\quad f\mapsto B\circ f
\end{equation*}
and $J:F^*\times F\rightarrow \Bbbk$ is the canonical pairing and $\Phi_J^p$ is the corresponding $L^p$-pairing. One 
can check this diagram commutes. Now it suffices to show that $\Phi_J^p$ is isometry, as if so, we will have:
\begin{equation*}
    \|\Phi_B^p (f)\|=\|\Phi_J^p\circ L^p(B)(f)\|=\|B\circ f\|=\|f\|
\end{equation*}
as both $\Phi_J^p$ and $B$ are isometries.

To show $\Phi_J^p$ is isometry, i.e. for all $h\in L^p(X;F^*)$, we have $\|\Phi_J^p(h)\|=\|h\|$. First 
consider the following lemma:
\begin{lemma}
For a continuous bilinear pairing $B:E\times F\rightarrow \Bbbk$ and its
corresponding $L^p$-pairing $\Phi_B^p$, we have $\|\Phi_B^p\|\leq \|B\|$.
\end{lemma}
\begin{proof}[Proof of Lemma]
Let $p,q$ be conjugates, we have
\begin{equation*}
    |\Phi_B^p(g)(f)|\leq \int_X |B(g(x),f(x))|d\mu(x)\leq \|B\|\int_X\|g(x)\|\|f(x)\| d\mu(x)\leq \|B\|\|g\|_p\|f\|_q
\end{equation*}
where the last step uses H\"{o}lder's inequality. Thus
\begin{equation*}
    \|\Phi_B^p\|=\sup_{\|f\|_q\leq 1, \|g\|_p\leq 1}\|\Phi_B^p(g)(f)\|\leq \|B\|.
\end{equation*}
\end{proof}
Therefore, given $h\in L^p(X;F^*)$, we have 
\begin{equation*}
    \|\Phi_J^p(h)\|\leq \|\Phi_J^p\|\|h\|\leq \|J\|\|h\|\leq \|h\|
\end{equation*}
as one can show that $\|J\|\leq 1$. So now it suffices to show that for given $h\in L^p(X;F^*)$, we 
have $\|\Phi_J^p(h)\|\geq \|h\|$. 

For $1\leq p <\infty$, we first consider the following lemma:
\begin{lemma}
Given $\varepsilon>0$, for all $\varphi\in F^*$, there exists $v\in F$ s.t. $\|v\|\leq 1$ and 
$\text{Re}\;\varphi(v)\geq (1-\varepsilon)\|\varphi\|$.
\end{lemma}
\begin{proof}[Proof of Lemma]
Note $\|\varphi\|=\sup_{\|v\|\leq 1}\|\varphi(v)\|$, so for all $\varepsilon>0$, we can always find $\|v\|\leq 1$ s.t.
\begin{equation*}
    |\varphi(v)|\geq (1-\varepsilon)\|\varphi\|
\end{equation*}
Now let's consider $\tilde{v}:=\frac{\overline{\varphi(v)}}{|\varphi(v)|}v$, then $\varphi(\tilde{v})=|\varphi(v)|$, so
\begin{equation*}
    \text{Re}\;\varphi(\tilde{v})=\text{Re}\;|\varphi(v)|=|\varphi(v)|>(1-\varepsilon)\|\varphi\|.
\end{equation*}
\end{proof}
Given $\varepsilon>0$, let's consider $h$ to be a simple function first and write $h=\sum_{i=1}^n \varphi_i \1_{A_i}$ with $A_i$ pairwise disjoint and $\varphi_i\neq 0$, so that $\mu(A_i)<\infty$.
Then for each $\varphi_i\in F^*$, pick $v_i\in F$ s.t. $\|v_i\|\leq 1$ and 
$\text{Re}\;\varphi_i(v_i)\geq (1-\varepsilon)\|\varphi_i\|$. Now consider 
\begin{equation*}
    f=\sum_{i=1}^n \|\varphi_i\|^{p-1}v_i \1_{A_i}
\end{equation*}
First, we can see that $f\in L^q$, as it is a simple function supported on a set of finite measure. Moreover,
\begin{equation*}
    \|f\|_q^q=\sum_{i=1}^n\|\varphi_i\|^{q(p-1)}\|v_i\|^q \mu(A_i)\leq \sum_{i=1}^n\|\varphi_i\|^p \mu(A_i)=\|h\|_p^p
\end{equation*}
thus $\|f\|_q \leq \|h\|_p^{p/q}=\|h\|_p^{p-1}$ (Here we skip the details for specific case of 
$p=1$, $q=\infty$, but it is not hard to check). Note that 
\begin{align*}
    |\Phi_J^p(h)(f)|&\geq \text{Re}\; \Phi_J^p(h)(f)\\
      &= \sum_i \|\varphi_i\|^{p-1}\text{Re}\;\varphi_i(v_i)\mu(A_i)\\
      &\geq (1-\varepsilon)\sum_i \|\varphi_i\|^p \mu(A_i)\\
      &= (1-\varepsilon)\|h\|_p^p.
\end{align*}
so we have
\begin{equation*}
    \|\Phi_J^p(h)\|\geq \frac{|\Phi_J^p(h)(f)|}{\|f\|_q}\geq (1-\varepsilon)\frac{\|h\|_p^p}{\|f\|_q}\geq (1-\varepsilon)\frac{\|h\|_p^p}{\|h\|_p^{p-1}}=(1-\varepsilon)\|h\|.
\end{equation*}
As it is true for all given $\varepsilon>0$, we have $\|\Phi_J^p(h)\|\geq \|h\|$ for $h$: simple function. Now let's consider
\begin{equation*}
    S:=\{g\in L^p(X;F^*):\|g\|\leq \|\Phi_J^p(g)\|\}\subset L^p(X;F^*).
\end{equation*}
It is easy to see that $S$ is closed and it contains all simple functions in $L^p$. As these simple functions are dense in 
$L^p$, $S=L^p$, which gives us desired result.

For $p=\infty$, assume $\mu$ is semifinite. We want to show that $\|\Phi_{J}^\infty(h)\|\geq \|h\|_\infty$ for 
$\Phi_{J}^\infty : L^\infty(X;F^*)\rightarrow (L^1(X;F))^*$. First we introduce the following lemma:
\begin{lemma}
For all measurable $D\subset X$ with $0<\mu(D)<\infty$, we have 
\begin{equation*}
    \|\Phi^\infty_J(h)\|\geq \frac{1}{\mu(D)}\int_D \|h(x)\| d\mu(x).
\end{equation*}
\end{lemma}
\begin{proof}[Proof of Lemma]
Given measurable $D\subset X$ with $0<\mu(D)<\infty$. Consider the operator 
\begin{equation*}
    \iota_D : L^\infty(X) \rightarrow L^1(X), \quad f\mapsto \1_D f
\end{equation*}
It is well-defined as $\mu(D)$ is finite and $f\in L^\infty$ is bounded a.e.. Then we can show that the following diagram
\begin{equation*}
\begin{tikzcd}
L^\infty (X;F^*) \arrow[rr, "\Phi_J^\infty"] \arrow[d, "\iota_D"'] &  & (L^1(X;F))^* \arrow[d, "\iota_D^*"] \\
L^1 (X;F^*) \arrow[rr, "\Phi_J^1"]                                 &  & (L^\infty(X;F))^*                  
\end{tikzcd}
\end{equation*}
commutes. This can be shown by:
\begin{align*}
  \Phi_J^1(\1_D h)(f)&=\int_X \1_D(x)h(x)(f(x))d\mu(x)\\
  &=\int_D h(x)(\1_D(x)f(x))d\mu(x)=\Phi^\infty_J(h)(\1_D f).
\end{align*}
Therefore,
\begin{equation*}
    |\Phi_J^1(\1_D h)(f)|=|\Phi^\infty_J(h)(\1_D f)|\leq \|\Phi^\infty_J(h)\|\|\1_D f\|_1\leq \mu(D)\|\Phi^\infty_J(h)\|\| f\|_\infty
\end{equation*}
This implies:
\begin{equation*}
    \|\Phi_J^1(\1_D h)\|=\sup_{\|f\|_\infty \leq 1}|\Phi_J^1(\1_D h)(f)|\leq \sup_{\|f\|_\infty \leq 1}\mu(D)\|\Phi^\infty_J(h)\|\| f\|_\infty\leq \mu(D)\|\Phi^\infty_J(h)\|.
\end{equation*}
Recall that we've already proved that 
\begin{equation*}
    \|\Phi_J^1(\1_D h)\|=\|\1_D h\|_1=\int_D \|h(x)\|d\mu(x)
\end{equation*}
Thus we have our desired result.
\end{proof}
Now let $M:=\|h\|_\infty$. For $M=0$, the inequality trivially holds. For $M>0$, given $0<\varepsilon<M$, consider
\begin{equation*}
    A:=\{x\in X: \|h(x)\|>M-\varepsilon\}
\end{equation*}
By definition of $\|\cd\|_\infty$, $\mu(A)>0$. As $\mu$ is semifinite, there exists $D\subset A$: measurable s.t. $0<\mu(D)<\infty$.
As $D\subset A$, for all $x\in D$, we have $\|h(x)\|>M-\varepsilon$. Therefore,
\begin{equation*}
    \|\Phi^\infty_J(h)\|\geq \frac{1}{\mu(D)}\int_D \|h(x)\|d\mu(x)\geq M-\varepsilon.
\end{equation*}
As it is true for all given $\varepsilon$, we have $\|\Phi_J^\infty(h)\|\geq \|h\|_\infty$.
\end{proof}


\begin{theorem}
Let $p$ and $q$ be H\"{o}lder conjugate, and let $B:\Bbbk\times\Bbbk\rightarrow\Bbbk$ be the multiplication.
If $1<p<\infty$, for each $\phi\in(L^q(X;\Bbbk))^*$ there exists $f\in L^p(X;\Bbbk)$ such that
$\phi(g)=\int_X fg\, d\mu$ for all $g\in L^q(X;\Bbbk)$, and hence
$\Phi_B^p:L^p(X;\Bbbk)\rightarrow(L^q(X;\Bbbk))^*$ is an isometric isomorphism.
The same conclusion holds for $p=\infty$ provided $\mu$ is $\sigma$-finite.
\end{theorem}















